% Problem record op_4ada32907dfeb9b5. This TeX file is derived from
% database/problems_json/op_4ada32907dfeb9b5.json by scripts/sync-tex.mjs; edit the JSON record
% and rerun the script rather than editing this file.
\section{GHZ distillation rate of a triangle of dephased Bell pairs}
\label{sec:op_4ada32907dfeb9b5}

\paragraph{Problem.}
Alice, Bob, and Charlie share $n$ copies of a triangle of dephased Bell pairs. Let $\Phi_2$ be the projector onto $(\ket{00}+\ket{11})/\sqrt{2}$ and $\Phi_2^-$ the projector onto $(\ket{00}-\ket{11})/\sqrt{2}$. An independent phase flip on each edge, with known probability $p$, gives the edge state $\sigma_p=(1-p)\Phi_2+p\Phi_2^-$. The state of one triangle is $\tau_p=\sigma_p^{AB}\otimes\sigma_p^{BC}\otimes\sigma_p^{CA}$. The parties receive only the edge qubits.

For a proposed rate $r$, the parties aim to distill $\lfloor rn\rfloor$ copies of the tripartite GHZ state with vanishing trace-distance error on the full target. Let $\mathrm{GHZ}_3$ be the projector onto $(\ket{000}+\ket{111})/\sqrt{2}$.

What is the optimal LOCC distillation rate $R_\triangle(p)$ for $0<p<1/2$? The rate is defined by
\begin{equation}
  R_\triangle(p)=\sup\left\{r\geq 0:\lim_{n\to\infty}\inf_{\Lambda_n\in\mathrm{LOCC}}
  \tfrac12\left\lVert\Lambda_n\bigl(\tau_p^{\otimes n}\bigr)
  -\mathrm{GHZ}_3^{\otimes\lfloor rn\rfloor}\right\rVert_1=0\right\}.
  \label{eq:deph-triangle-rate}
\end{equation}
In Eq.~\eqref{eq:deph-triangle-rate}, the infimum is over LOCC channels with $\lfloor rn\rfloor$ output qubits per party. A local phase flip interchanges $p$ and $1-p$ without changing the rate.

Give an achievable rate and a matching converse for arbitrary LOCC for $0<p<1/2$, counting GHZ states per triangle.

\subsection*{Status}

Unsolved

\subsection*{Source}

This question asks how joint processing of a triangle of independently dephased Bell pairs compares with distilling each link separately. The single-link hashing protocol of Bennett et al. gives a starting point \sourcecite{ref:feb9b5-bdsw96}{BDSW96} (Section 3.2.3).

\subsection*{Progress}

\begin{itemize}
  \item Let $d(p)=1-h_2(p)$, where $h_2$ is the binary entropy in bits. This is the bipartite distillable entanglement of $\sigma_p$. One-way hashing achieves it \sourcecite{ref:feb9b5-bdsw96}{BDSW96} (Section 3.2.3), and Rains gives the matching converse for maximally correlated states \sourcecite{ref:feb9b5-rai01}{Rai01} (Theorem 6.2). The bounds are given in Eq.~\eqref{eq:deph-triangle-bounds}.
\begin{equation}
  \tfrac32 d(p)\leq R_\triangle(p)\leq 2d(p).
  \label{eq:deph-triangle-bounds}
\end{equation}
For the lower bound, distill Bell pairs independently on each edge and use the three possible two-edge spanning trees equally often. Two Bell pairs meeting at a vertex give one GHZ state.

For the upper bound, group Bob and Charlie into one laboratory. Two links cross this cut per triangle, giving $2d(p)$ distillable ebits. Each GHZ state carries one ebit across it.
\end{itemize}

\subsection*{References}

\begin{enumerate}[label={},leftmargin=4.8em,labelsep=0.5em]
  \footnotesize
  \item[\textup{[BDSW96]}]\label{ref:feb9b5-bdsw96}
  C. H. Bennett, D. P. DiVincenzo, J. A. Smolin, W. K. Wootters, “Mixed State Entanglement and Quantum Error Correction”, 1996.
\href{https://arxiv.org/abs/quant-ph/9604024}{arXiv:quant-ph/9604024}

  \item[\textup{[Rai01]}]\label{ref:feb9b5-rai01}
  E. M. Rains, “A semidefinite program for distillable entanglement”, 2001.
\href{https://arxiv.org/abs/quant-ph/0008047}{arXiv:quant-ph/0008047}
\end{enumerate}

\subsection*{Comment}

At $p=0$ the links are perfect and the spanning-tree protocol achieves $3/2$ GHZ states per triangle. At $p=1/2$ every link is separable and $R_\triangle(1/2)=0$.

\subsection*{Field}

Quantum Resource Theory

\subsection*{Topic}

Entanglement cost; Local operations and classical communication; Entanglement distillation

\subsection*{ID}

\texttt{op\_4ada32907dfeb9b5}
